\chapter{The expression language}\label{ch:expr}
\begin{goals}
\textbf{In this chapter} you get the full language lamina's filters, calculated columns and grade–tonnage
factors are written in: names, values, operators, the rules for blank values, and every function.
\end{goals}

One language does three jobs in lamina. In the filter box an expression is a condition, true or false
for each row. In a calculated column it is a value: a number or a text. In grade–tonnage it gives the
volume, density or ore proportion of each block. The language is the same in all three, and it is the
same one micro uses.

It looks like the \code{WHERE} clause of SQL, because that is what people who work with tables already
read:

\begin{codeblock}
LITO in ("SULPHIDE", "TRANSITION") and Cu_pct between 0.5 and 2
\end{codeblock}

Two properties matter before any detail. An expression is always about one row: it can use any column
of that row, and nothing outside it (no previous row, no total). And an expression never fails while it
runs: arithmetic that has no answer (a division by zero, the square root of a negative number) gives a
blank value, not an error, so a filter over ten million rows can't stop at row six million because one
value was odd (\cref{sec:totality}).

\section{Names and values}\label{sec:names}

\subsection{Columns}

A word is a column name: \code{Cu\_pct}, \code{LITO}, \code{dx}. Names are matched without regard to
case, so \code{cu\_pct} and \code{CU\_PCT} are the same column. A plain name is letters, digits and
underscores, not starting with a digit.

Any other name goes between backticks:

\begin{codeblock}
`Cu (ppm)` > 30
`OK-Indic` = 1
`in` = "yes"
\end{codeblock}

That covers spaces, brackets, hyphens and names that are also keywords. Two other forms are read too:
\code{\$\{site\_id\}}, the reference form of XLSForm, for names made of letters, digits, \code{\_},
\code{.} and \code{-} (so \code{\$\{1st\_pass\}} and \code{\$\{fe-pct\}} work); and \code{["Cu (ppm)"]},
an older form lamina's lenses may still hold.

\begin{caution}
A hyphen between two words is always a minus: \code{OK-Indic} is the column \col{OK} minus the column
\col{Indic}. When a column really is called \col{OK-Indic}, lamina's error says so and suggests the
backticked form.
\end{caution}

\subsection{Values}

\begin{center}\small
\begin{tabularx}{\linewidth}{@{}lL@{}}
\toprule
\sffamily\bfseries write & \sffamily\bfseries for \\
\midrule
\code{42}, \code{3.14}, \code{.5}, \code{1.5e-3} & numbers \\
\code{"OXIDE"}, \code{'OXIDE'} & text (double quotes are the usual form) \\
\code{true}, \code{false} & the two truth values \\
\code{blank} & a missing value (\cref{sec:blanks}) \\
\bottomrule
\end{tabularx}
\end{center}

A word without quotes is a column; a word in quotes is text. So \code{fe > cu} compares two columns, and
\code{LITO = "OXIDE"} compares a column with a text. This is the one rule that trips people coming from a
spreadsheet: \code{LITO = OXIDE} looks for a column called OXIDE.

\subsection{Text that looks like a number}

A text value that reads as a number compares as a number: \code{"10" < "9"} is false, and a column
whose cells hold \code{12} as text equals \code{12}. Text that isn't a number compares as text, letter by
letter, with capitals before small letters (\code{"abc" < "abd"} is true). A file read with a decimal
comma (\cref{sec:decimal-comma}) reads \code{1,5} in its cells as 1.5 here too.

\section{Operators}\label{sec:operators}

\begin{center}\small
\begin{tabularx}{\linewidth}{@{}lL@{}}
\toprule
\sffamily\bfseries operator & \sffamily\bfseries meaning \\
\midrule
\code{+ - * /} & arithmetic; unary minus \\
\code{\textasciicircum{}} or \code{**} & power; \code{2\textasciicircum{}3\textasciicircum{}2} is $2^{(3^2)} = 512$ and \code{-2\textasciicircum{}2} is $-4$ \\
\code{= != < <= > >=} & comparison \\
\code{between} … \code{and} … & a range, ends included: \code{Z between 400 and 500} \\
\code{in (…)}, \code{not in (…)} & one of a set: \code{LITO in ("OXIDE", "SULPHIDE")} \\
\code{contains}, \code{not contains} & text containing other text \\
\code{like}, \code{not like} & text matching a pattern with \code{\%} and \code{\_} \\
\code{matches} & text matching a regular expression \\
\code{is blank}, \code{is filled}, \code{is not blank} & whether a value is missing \\
\code{and}, \code{or}, \code{not} & combining conditions \\
\code{( )} & grouping \\
\bottomrule
\end{tabularx}
\end{center}

From tightest to loosest: the power; unary minus; \code{* /}; \code{+ -}; the comparisons and the
\code{between}, \code{in}, \code{contains}, \code{like}, \code{matches} and \code{is} tests; \code{not};
\code{and}; \code{or}. So \code{a or b and c} is \code{a or (b and c)}, \code{not Z > 500} is
\code{not (Z > 500)}, and \code{2 * 3\textasciicircum{}2} is 18. Comparisons don't chain:
\code{400 < Z < 500} is an error; write \code{Z between 400 and 500}. When in doubt, add brackets: they
cost nothing.

\subsection{Other spellings}

The C and spreadsheet spellings are read too, so habits from elsewhere work: \code{==} for \code{=},
\code{<>} for \code{!=}, \code{\&\&} and \code{||} for \code{and} and \code{or}, \code{\textasciitilde{}} and
\code{!\textasciitilde{}} for \code{contains} and \code{not contains}, single quotes for text, and \code{in "a", "b"}
without brackets. lamina itself writes \code{==} when it builds a filter from a cell
(\cref{sec:filter-from-grid}). Words (\code{and}, \code{between}, \code{blank}) may be in any case.

\subsection{Matching text}\label{sec:text-match}

\code{contains}, \code{like} and \code{matches} are case-sensitive: \code{LITO contains "XID"} is true
for OXIDE, \code{LITO contains "xid"} isn't. To ignore case, change it first:
\code{upper(LITO) contains "XID"}.

\code{like} takes \code{\%} for any run of characters and \code{\_} for one: \code{HOLE like "DDH\%"},
\code{HOLE like "DDH-01\_"}. \code{matches} takes a regular expression: \code{HOLE matches "\textasciicircum{}DDH-0[0-9]+\$"}.
There is no \code{not matches}; write \code{not (HOLE matches "…")}.

A regular expression that could take an astronomical time on some text is refused when you type it,
rather than allowed to freeze the window. These are patterns with a repeated group that itself
contains a repetition or an alternative: \code{(a+)+}, \code{(x|xx)+}, \code{(a|b)*}, and
back-references. Rewrite without the nesting: \code{[ab]*} for \code{(a|b)*}.

\section{Blank values}\label{sec:blanks}

A cell can be empty, and many files mark a missing value some other way that lamina reads as empty: a
Datamine absent value (\cref{sec:dm}), a missing Parquet value. In an expression these are
\term{blank}. How blanks behave is the most important rule in the language, because a missing grade
that quietly became zero would pull every mean and every tonnage down with it.

\begin{description}
\item[Blank goes through arithmetic.] \code{Cu\_pct * 2} is blank where \col{Cu\_pct} is blank, and so
is \code{Cu\_pct + 1} and \code{Cu\_pct * 0}. A blank never turns into zero by itself.
\item[You turn it into a value on purpose.] \code{coalesce(Cu\_pct, 0)} is \col{Cu\_pct}, or 0 where it
is blank; \code{ifnum(x, 0)} is \code{x} where \code{x} is a number and 0 otherwise. These, with
\code{if}, are the only ways a blank becomes a value.
\item[A blank never satisfies an order.] \code{Cu\_pct > 1} and \code{Cu\_pct <= 1} are both false where
\col{Cu\_pct} is blank: a missing grade is neither high nor low.
\item[Equality is friendly.] \code{x = blank} is true where \code{x} is blank, \code{blank = blank} is
true, and \code{Cu\_pct = 0} is false and \code{Cu\_pct != 0} true where it is blank. There is no SQL
trap where \code{= NULL} matches nothing. \code{x is blank} is the clearer spelling.
\item[\code{not} is the complement.] \code{not (Cu\_pct > 1)} is true where \col{Cu\_pct} is blank. A
filter and \code{not} of it always split the rows between them, blanks included; in SQL the blanks fall
out of both.
\end{description}

\begin{center}\small
\begin{tabularx}{\linewidth}{@{}>{\ttfamily}lcL@{}}
\toprule
\normalfont\sffamily\bfseries where Cu\_pct is blank & \sffamily\bfseries gives & \\
\midrule
Cu\_pct > 1 & false & a blank satisfies no order \\
not (Cu\_pct > 1) & true & the complement takes it \\
Cu\_pct <= 1 & false & so \code{not (x > 1)} is not \code{x <= 1} on blanks \\
Cu\_pct = blank & true & \\
Cu\_pct != 0 & true & \\
Cu\_pct * 2 & blank & \\
coalesce(Cu\_pct, 0) & 0 & \\
if(Cu\_pct > 1, 1, 0) & 0 & a condition that isn't true takes the \code{else} \\
\bottomrule
\end{tabularx}
\end{center}

\code{and} and \code{or} are strict about truth: they take \code{true} as true and everything else,
blanks and numbers included, as false. \code{Cu\_pct and true} is false even where \col{Cu\_pct} is 2.5.
Write the comparison you mean: \code{Cu\_pct > 0 and …}.

An empty text and a blank are the same thing: a function whose result is an empty text gives blank.

\section{Functions}\label{sec:functions}

Every function gives a blank when an argument it needs is blank or makes no sense, rather than an error.
\Cref{app:functions} lists them all alphabetically; here they are by use.

\subsection{Choosing}

\begin{description}
\item[\code{if(cond, then, else)}] \code{then} where \code{cond} is true, \code{else} otherwise:
\code{if(Cu\_pct >= 0.3, "ore", "waste")}. With \code{blank} as a branch it scrubs a value:
\code{if(AU = -99, blank, AU)}.
\item[\code{coalesce(a, b, …)}] the first that isn't blank.
\item[\code{ifnum(x, default)}] \code{x} if it is a number, else \code{default}.
\item[\code{nullif(x, v)}] blank where \code{x} equals \code{v}: \code{nullif(AU, -999)} removes the
$-999$ missing-value code.
\end{description}

\subsection{Numbers}

\begin{description}
\item[\code{round(x)}, \code{round(x, n)}] to the nearest whole number, or to \code{n} decimals; a half
rounds up, towards $+\infty$ (\code{round(2.5)} is 3, \code{round(-2.5)} is $-2$). Numbers are binary
floating point, so a value like 1.005 is stored a hair below itself and \code{round(1.005, 2)} is 1.
\item[\code{int(x)}] drops the fraction, towards zero: \code{int(-2.7)} is $-2$.
\item[\code{floor(x)}, \code{ceil(x)}] down and up to a whole number.
\item[\code{abs(x)}] the absolute value.
\item[\code{mod(x, y)}] the remainder, with the sign of \code{y}, as in Excel: \code{mod(-7, 3)} is 2;
\code{mod(x, 0)} is blank.
\item[\code{sqrt}, \code{exp}, \code{log}, \code{log10}, \code{pow(x, y)}] \code{log} is the natural
logarithm, \code{log10} the base-10 one geochemists log grades with. \code{log(0)}, \code{sqrt(-1)} and
the like are blank.
\item[\code{min(…)}, \code{max(…)}] the least and greatest of their arguments, leaving blanks out:
\code{max(3, Cu\_pct)} is 3 where \col{Cu\_pct} is blank.
\item[\code{clamp(x, lo, hi)}] \code{x}, held between \code{lo} and \code{hi}.
\item[\code{bin(x, w)}, \code{bin(x, w, origin)}] the lower edge of the band of width \code{w} that
holds \code{x}: \code{bin(7.3, 2)} is 6, \code{bin(-1, 2)} is $-2$, and with an origin,
\code{bin(7.3, 2, 1)} is 7. As a calculated column it gives benches, elevation bands and grade classes to
group by (\cref{sec:bands}).
\end{description}

\subsection{Interpolation}

\begin{description}
\item[\code{lerp(a, b, t)}] the value a fraction \code{t} of the way from \code{a} to \code{b}.
\item[\code{remap(x, a1, b1, a2, b2)}] \code{x} taken from the range \code{a1}–\code{b1} to the range
\code{a2}–\code{b2}: \code{remap(5, 0, 10, 100, 200)} is 150. Blank when \code{a1} equals \code{b1}.
\item[\code{interp(x, x1, y1, x2, y2, …)}] a piecewise-linear curve through the points, flat beyond the
first and last: \code{interp(Fe\_pct, 20, 2.7, 45, 3.4, 65, 4.3)} is a density from an iron grade,
without a ladder of \code{if}s. Written as numbers, the \code{x}s must increase, or the expression is
refused as a likely typo; the points are a list of pairs, so the count of arguments after \code{x} must
be even.
\item[\code{smoothstep(a, b, x)}] 0 below \code{a}, 1 above \code{b}, an S-curve between.
\end{description}

\code{lerp} and \code{remap} don't hold their result inside the range; wrap them in \code{clamp} for
that.

\subsection{Text}

\begin{description}
\item[\code{upper}, \code{lower}, \code{trim}] case, and spaces at either end.
\item[\code{len(s)}] the number of characters; blank for a blank.
\item[\code{left(s, n)}, \code{right(s, n)}] the first or last \code{n} characters.
\item[\code{substr(s, start)}, \code{substr(s, start, n)}] from position \code{start}, counting from 1:
\code{substr("DDH-012", 5, 3)} is \code{012}.
\item[\code{replace(s, find, with)}] every occurrence, taken literally.
\item[\code{concat(…)}] joined, leaving blanks out, numbers written as text:
\code{concat("DDH-", 23)} is \code{DDH-23}.
\end{description}

\code{upper(trim(HOLEID))} is the usual way to make hole names written by different people agree.

\subsection{Dates}

\code{year(s)}, \code{month(s)}, \code{day(s)} take a date written \code{YYYY-MM-DD} (as text; lamina has
no date type) and give its parts.

\subsection{Tests}

\code{isnum(x)} is true where \code{x} reads as a number (\code{"1e3"} and \code{" 12 "} do);
\code{isnan(x)} where it is present but isn't a number; \code{isblank(x)} and \code{isfilled(x)} where
it is missing or not. As filters they find the junk in a numeric column: \code{isnan(Cu\_pct)} lists the
rows where \col{Cu\_pct} holds something like \code{<0.01} or \code{n/a}.

\section{Long expressions}\label{sec:comments}

An expression can run over several lines, and \code{\#} starts a comment that runs to the end of its
line. In the filter box and the calculated-column editor, that lets a long condition be laid out and
explained:

\begin{codeblock}
LITO in ("SULPHIDE", "TRANSITION")   # the fresh rock
and Cu_pct >= 0.3                    # the reporting cutoff
and Z between 400 and 520            # above the 2030 pit floor
\end{codeblock}

\section{It never fails as it runs}\label{sec:totality}

An expression that is read correctly runs to the end on every row. Division by zero, the logarithm of a
negative, text where a number was wanted: each gives a blank for that row, and the scan goes on. Mistakes
are caught when you type the expression, not millions of rows into a scan: an unknown column, a
missing bracket, a function given the wrong number of arguments (\code{round() expects 1–2 argument(s),
got 0}), an \code{interp} whose points are out of order, a dangerous regular expression.

\begin{why}
This is what makes an expression safe to receive from someone else: in a lens, from a colleague, from a
script. It can't loop, can't stop with an error halfway, can't reach anything outside the row, and lamina
runs it without turning it into program code (\cref{sec:no-eval}). The worst a strange expression can do
is give blanks.
\end{why}

\begin{tryit}
Blanks, on purpose.
\begin{enumerate}
\item In the sample, add a calculated column (\ui{ƒ+ col}) called \col{cu\_hi} with the expression
\code{if(Cu\_pct >= 1, Cu\_pct, blank)}.
\item Filter \code{cu\_hi is blank}. These are the blocks under 1\%: most of the model.
\item Change the filter to \code{not (cu\_hi > 2)}. The blocks between 1 and 2\% and every blank: the
complement of \code{cu\_hi > 2}.
\item Clear the filter. Right-click the \col{cu\_hi} header: if it offers \ui{Treat as number}, choose it
(lamina guessed the new column's type from rows that were all blank, \cref{sec:calc-type}). Then choose
\ui{Statistics — cu\_hi…}: the blanks are counted as blanks, and the mean is the mean of the blocks at or
above 1\% only.
\end{enumerate}
\end{tryit}
